\documentclass[11pt,a4paper]{article}
\usepackage[margin=25mm]{geometry}
\usepackage{fontspec}
\setmainfont{TeX Gyre Pagella}
\usepackage{amsmath,amssymb,amsthm,mathtools}
\usepackage{enumitem,microtype,xcolor}
\usepackage[colorlinks=true,linkcolor=blue!45!black,
  citecolor=blue!45!black,urlcolor=blue!45!black]{hyperref}
\setlength{\emergencystretch}{3em}
\setlength{\parskip}{3pt}
\setlist{itemsep=3pt,topsep=5pt}
\allowdisplaybreaks[2]
\newcommand{\C}{\mathbb C}
\newcommand{\Z}{\mathbb Z}
\newcommand{\N}{\mathbb N}
\newcommand{\g}{\mathfrak g}
\newcommand{\slc}{\mathfrak{sl}_2}
\newcommand{\ad}{\operatorname{ad}}
\newcommand{\Span}{\operatorname{span}}
\newcommand{\T}{\mathcal T}
\newcommand{\I}{\mathcal I}
\newtheoremstyle{plainupright}{6pt}{5pt}{\normalfont}{}{\bfseries}{.}{.5em}{}
\theoremstyle{plainupright}
\newtheorem{theorem}{Theorem}[section]
\newtheorem{proposition}[theorem]{Proposition}
\newtheorem{lemma}[theorem]{Lemma}
\newtheorem{corollary}[theorem]{Corollary}
\newtheorem{remark}[theorem]{Remark}
\hypersetup{
  pdftitle={Additional Progress on Localization Quotients},
  pdfsubject={Deep-region simplicity, the unique maximal submodule,
    and compatibility with central reduction}}
\title{Additional Progress on Localization Quotients\\[4pt]
  \large Deep-region simplicity, the unique maximal submodule,\\
  and compatibility with central reduction}
\author{}
\date{September 15, 2026}
\begin{document}
\maketitle
\begin{abstract}
This addendum completes the nonintegral simplicity arguments for
\(\T_{f_0}M_{-1,1}\) and, more generally, for the deep-region quotients
\(\T_{f_c}M_{a,b}\) of simple induced modules.
The main new argument shows that every proper submodule of such a quotient
is contained in its \(e_{-c}\)-locally nilpotent part.
At integral parameters this part is therefore the unique maximal submodule;
the corresponding quotient is simple, the extension is nonsplit, and the
localization quotient is indecomposable.
We also prove that algebraic Sugawara central reduction commutes with
quotient localization, including at nonboundary modes.
Simplicity of the maximal submodule itself and shallow-region simplicity
are not established here.
\end{abstract}

\section{Definitions, assumptions, and the new statements}

We work over \(\C\), write \(\N=\{1,2,\ldots\}\), and initially omit
the degree derivation. Set
\[
\g=(\slc\otimes\C[t,t^{-1}])\oplus\C K,\qquad U=U(\g),
\]
where \(x_r=x\otimes t^r\), with brackets
\begin{align*}
[e_r,f_s]&=h_{r+s}+r\delta_{r+s,0}K,&
[h_r,h_s]&=2r\delta_{r+s,0}K,\\
[h_r,e_s]&=2e_{r+s},&
[h_r,f_s]&=-2f_{r+s}.
\end{align*}
The \(e\)'s commute, the \(f\)'s commute, and \(K\) is central.

Fix \(a,b\in\Z\) with \(a+b\ge0\), put \(L=a+b+1\ge1\), and define
\[
S_{a,b}=\Span\{K,\ h_i\ (i\ge0),\ e_i\ (i>a),\ f_j\ (j>b)\}.
\]
Let \(\chi:S_{a,b}\to\C\) be the character with
\[
\chi(K)=\kappa,\qquad \chi(h_i)=\lambda_i,\qquad
\chi(e_i)=\chi(f_j)=0,\qquad \lambda_i=0\quad(i>L).
\]
Let \(\C_\chi\) denote the corresponding one-dimensional module, and set
\[
M=M_{a,b}(\chi)=U\otimes_{U(S_{a,b})}\C_\chi,\qquad u=1\otimes1_\chi.
\]
If the notation \(M^c_{a,b}\) from the unified note \cite{Unified} is used,
it means this same module \(M_{a,b}\), without the derivation.

For a root mode \(F=f_c\), the powers of \(F\) form an Ore set:
\(\ad F\) is locally nilpotent on \(U\).
Write
\[
U_F=U[F^{-1}],\qquad D_FV=U_F\otimes_U V,\qquad
\T_FV=\operatorname{coker}(V\longrightarrow D_FV).
\]
When \(F\) is injective on \(V\), we identify \(V\) with its image and
write \(\T_FV=D_FV/V\). Negative powers act on \(D_FV\), not on
\(\T_FV\).

For the simplicity results in Sections~2--4, our affine hypotheses are
\begin{equation}\label{eq:affine-assumptions}
\lambda_L\ne0,\qquad \kappa\ne-2,\qquad
c<b,\qquad c\le-a-1.
\end{equation}
The last condition is the \emph{deep-region} condition
\[
d_*:=b-c\ge L.
\]
Put
\[
E=e_{-c},\qquad F=f_c,\qquad H=h_0-cK,\qquad
\mu=\lambda_0-c\kappa.
\]
Then
\[
[H,E]=2E,\qquad [H,F]=-2F,\qquad [E,F]=H,\qquad
Eu=0,\qquad Hu=\mu u.
\]
By the FGXZ simplicity criterion, the first two conditions in
\eqref{eq:affine-assumptions} make \(M\) simple
\cite[Theorem 1.1]{FGXZ}.

Order PBW with \(F\) first, followed by the remaining free \(f\)-modes,
the free \(e\)-modes, the negative \(h\)-modes, and finally \(S_{a,b}\).
PBW then makes \(M\) free over \(\C[F]\), so \(F\) is injective. Define
\[
D=D_FM,\qquad Q_c=D/M,\qquad
w_m=F^{-m}u+M\quad(m\ge1).
\]
The localized PBW basis shows that the \(w_m\) are nonzero and that
\(F\) on \(Q_c\) is surjective and locally nilpotent. The operator \(H\)
is semisimple on \(M,D,Q_c\), with weights in \(\mu+2\Z\).
An operator is \emph{locally nilpotent} if every vector is killed by
some power of it; it is \emph{locally finite} if every vector generates
a finite-dimensional module for the polynomial algebra in that operator.

Set
\[
G_c=\{v\in Q_c:E^nv=0\text{ for some }n\ge1\}.
\]
The new conclusions will be
\[
Q_c\text{ is simple}\quad\Longleftrightarrow\quad\mu\notin\Z,
\]
and, when \(\mu\in\Z\), that \(G_c\) is the unique maximal submodule,
\(Q_c/G_c\) is simple, and their extension is nonsplit.
The argument is first proved in a general rank-one setting.

\section{A rank-one submodule argument}

\subsection{General hypotheses and Casimir local finiteness}

In this section let \(\mathfrak a\) be a complex Lie algebra containing
an \(\slc\)-triple \((E,H,F)\), with the brackets displayed above.
Write \(A=U(\mathfrak a)\), and assume:
\begin{enumerate}[label=\textup{(\roman*)}]
\item \(\ad E\) and \(\ad F\) are locally nilpotent on \(A\), and
the Ore localization \(A_F\) at the powers of \(F\) exists;
\item \(M\) is a simple \(A\)-module on which \(H\) acts semisimply,
\(E\) acts locally nilpotently, and \(F\) acts injectively.
\end{enumerate}
Let
\[
D=A_F\otimes_A M,\qquad Q=D/M,\qquad
G=\{q\in Q:E^nq=0\text{ for some }n\ge1\}.
\]
The \(H\)-action on \(D\) and \(Q\) is semisimple because
\([H,F^{-m}]=2mF^{-m}\).
The subspace \(G\) is an \(A\)-submodule: if \(E^nq=0\), the finite
Leibniz expansion
\[
E^r(sq)=\sum_{j=0}^r\binom rj(\ad E)^j(s)E^{r-j}q
\]
vanishes for sufficiently large \(r\), for every \(s\in A\).

Define the rank-one Casimir
\begin{equation}\label{eq:rank-one-C}
C=H^2+2H+4FE.
\end{equation}
It commutes with \(E,H,F\). We do \emph{not} assume that it is central
in \(A\). In particular, it is not the completed affine Casimir
\(\Omega\) involving the degree derivation.

\begin{lemma}[Casimir factorization]\label{lem:casimir}
For \(Hz=\nu z\) and \(n\ge0\),
\begin{equation}\label{eq:factorization}
\begin{aligned}
F^nE^nz&=A_n(C;\nu)z,\\
A_n(t;\nu)&=4^{-n}\prod_{j=0}^{n-1}
\bigl(t-(\nu+2j)(\nu+2j+2)\bigr),
\end{aligned}
\end{equation}
where \(A_0=1\). Consequently, \(C\) acts locally finitely on \(M\)
and on \(D\).
\end{lemma}
\begin{proof}
For \(n=1\), the identity is \(FE=(C-H^2-2H)/4\).
If it holds for \(n\), apply it to the weight vector \(Ez\), whose
weight is \(\nu+2\), and then multiply by \(F\). Since \(C\) commutes
with \(F\) and \(E\), this gives the next factor and proves
\eqref{eq:factorization}.

For a weight vector \(z\in M\), choose \(n\ge1\) with \(E^nz=0\).
Then the nonzero polynomial \(A_n(t;\nu)\) annihilates \(z\) when
evaluated at \(C\). Every vector of \(M\) has finitely many weight
components, so a product of such polynomials annihilates it.
Thus \(C\) acts locally finitely on \(M\).
Every vector of \(D\) can be written \(F^{-m}z\), and \(C\) commutes
with \(F^{-1}\). The same annihilating polynomial works in \(D\).
\end{proof}

\subsection{The infinitely divisible subspace}

\begin{lemma}\label{lem:divisible}
For an \(A\)-submodule \(P\subseteq D\), define
\[
W(P)=\bigcap_{n\ge0}F^nP.
\]
Then \(W(P)\) is an \(A\)-submodule and is stable under \(F^{-1}\).
If \(W(P)\ne0\), then \(P=D\).
\end{lemma}
\begin{proof}
Fix \(s\in A\), and choose \(\ell\ge0\) with
\((\ad F)^{\ell+1}(s)=0\). For every \(n\ge0\),
\[
sF^{n+\ell}
=\sum_{j=0}^{\ell}(-1)^j\binom{n+\ell}{j}
F^{n+\ell-j}(\ad F)^j(s).
\]
If \(w\in W(P)\), write \(w=F^{n+\ell}p\) with \(p\in P\).
Every term in the displayed expansion belongs to \(F^nP\), since
each \((\ad F)^j(s)\) preserves \(P\).
Hence \(sw\in F^nP\) for every \(n\), proving \(A\)-stability.
Notice that the individual spaces \(F^nP\) need not be \(A\)-submodules.

For \(w\in W(P)\), write \(w=F^{n+1}p\) for each \(n\ge0\).
Since \(F\) is invertible on \(D\), this gives \(F^{-1}w=F^np\).
Thus \(F^{-1}w\in W(P)\).

If \(0\ne w\in W(P)\), multiplication by a sufficiently large power
of \(F\) yields a nonzero vector of \(M\). It follows that
\(W(P)\cap M\ne0\). Simplicity of \(M\) gives \(M\subseteq W(P)\);
stability under \(F^{-1}\) then gives
\(D=D_FM\subseteq W(P)\subseteq P\subseteq D\).
\end{proof}

\begin{theorem}[Containment of all proper submodules]\label{thm:containment}
Under the general hypotheses of this section, every proper
\(A\)-submodule \(N\subsetneq Q\) is contained in \(G\).
If \(G\ne Q\), then \(G\) is the unique maximal submodule of \(Q\),
the quotient \(Q/G\) is simple, and \(Q\) is indecomposable.
If moreover \(G\ne0\), the sequence
\[
0\longrightarrow G\longrightarrow Q\longrightarrow Q/G
\longrightarrow0
\]
is nonsplit.
\end{theorem}
\begin{proof}
Let \(N\subseteq Q\) be a submodule not contained in \(G\), and let
\(P\subseteq D\) be its inverse image, so \(M\subseteq P\) and \(P/M=N\).
We show \(N=Q\).

Choose \(q\in N\) on which \(E\) is not locally nilpotent.
Since \(H\) is semisimple and every vector has finitely many weight
components, one weight component also has this property.
Each component belongs to \(N\), being obtainable by a polynomial
in \(H\). We may therefore assume that \(q\) is a weight vector.
Lift it to a weight vector \(w\in P\), say \(Hw=\eta w\).
By Lemma~\ref{lem:casimir}, choose a nonzero polynomial \(p(t)\)
such that \(p(C)w=0\).

There are only finitely many nonnegative integers \(k\) for which
\[
p\bigl((\eta+2k)(\eta+2k+2)\bigr)=0:
\]
each root of \(p\) gives at most two possible complex values of \(k\).
Choose an integer \(r\ge0\) larger than all such nonnegative integers,
and set
\[
v=E^rw,\qquad \nu=\eta+2r.
\]
Its image \(E^rq\) in \(Q\) is nonzero, so \(v\ne0\).
Also \(Hv=\nu v\) and \(p(C)v=0\), since \(C\) commutes with \(E\).
By the choice of \(r\),
\[
p\bigl((\nu+2j)(\nu+2j+2)\bigr)\ne0\qquad(j\ge0).
\]

For every \(n\ge1\), the polynomials \(p(t)\) and \(A_n(t;\nu)\)
are therefore relatively prime. Choose \(B_n(t),R_n(t)\in\C[t]\)
with
\[
A_n(t;\nu)B_n(t)+p(t)R_n(t)=1.
\]
The vector \(B_n(C)v\) has weight \(\nu\). Applying
\eqref{eq:factorization} to it gives
\[
F^n\bigl(E^nB_n(C)v\bigr)
=A_n(C;\nu)B_n(C)v=v.
\]
Since \(P\) is an \(A\)-submodule, \(E^nB_n(C)v\in P\).
Thus \(0\ne v\in W(P)\), and Lemma~\ref{lem:divisible} gives
\(P=D\), hence \(N=Q\).
This proves the containment assertion.

If \(G\ne Q\), it is itself a proper submodule and contains every
proper submodule. Hence it is the unique maximal submodule, and
\(Q/G\) is simple.
In a nontrivial direct sum decomposition of \(Q\), both summands
would be proper and therefore contained in \(G\), an impossibility.
This proves indecomposability.
When \(0\ne G\ne Q\), a splitting of the displayed exact sequence
would give such a nontrivial direct sum decomposition.
\end{proof}

\begin{remark}
The proof starts with an arbitrary submodule vector and constructs an
actual nonzero \(A\)-submodule \(W(P)\) inside the full localization.
This is the step not supplied by cyclicity of a chosen pure fraction.
No equivalence of BGG categories, finite weight multiplicity, or
centrality of \(C\) in \(A\) is used.
\end{remark}

\section{Application to the deep-region affine quotients}

Return to the affine notation of Section~1.
Both \(\ad E\) and \(\ad F\) are locally nilpotent on \(U\), as follows
directly from the current brackets and the Leibniz rule.
Since \(Eu=0\),
\[
E^n(su)=(\ad E)^n(s)u\qquad(s\in U),
\]
so \(E\) is locally nilpotent on the entire original module \(M\).
Thus all hypotheses of Section~2 hold under
\eqref{eq:affine-assumptions}.

For completeness, we include the rank-one calculation from the
unified note that determines when \(G_c\) is zero or proper.
This calculation itself does not require simplicity of \(M\).

\begin{proposition}[The locally nilpotent part]\label{prop:G}
Assume only \(c<b\) and \(c\le-a-1\), with the induced module and
character defined in Section~1. Then \(G_c\) is always a proper
submodule of \(Q_c\), and
\[
G_c=0\quad\Longleftrightarrow\quad \mu\notin\Z.
\]
If \(\mu\in\Z\) and \(k\) is any integer with
\(k\ge\max(1,\mu+2)\), then
\[
0\ne f_b^kw_1\in G_c,\qquad
E^{\,2k-\mu-1}(f_b^kw_1)=0.
\]
\end{proposition}
\begin{proof}
The rank-one commutation identity gives
\begin{equation}\label{eq:pure-recursion}
Ew_m=-m(\mu+m+1)w_{m+1},\qquad Hw_m=(\mu+2m)w_m.
\end{equation}
Choose a positive integer \(m\) such that
\(\mu+m+j+1\ne0\) for every \(j\ge0\).
Then all successive \(E\)-images of \(w_m\) are nonzero by PBW.
Hence \(w_m\notin G_c\), so \(G_c\ne Q_c\).

Suppose \(\mu\notin\Z\) and \(G_c\ne0\).
Taking a weight component and then its last nonzero \(E\)-image
produces a nonzero weight vector \(v\in\ker E\cap G_c\).
Write \(Hv=\nu v\). Since \(F\) is locally nilpotent on \(Q_c\),
choose the least \(n\ge1\) with \(F^nv=0\).
The rank-one identity yields
\[
0=EF^nv=n(\nu-n+1)F^{n-1}v.
\]
Minimality implies \(\nu=n-1\in\Z\), contradicting
\(\nu\in\mu+2\Z\). Thus \(G_c=0\).

Now let \(\mu\in\Z\), fix \(k\ge\max(1,\mu+2)\), and put
\[
Y=f_b,\qquad \eta=\lambda_{d_*}.
\]
On \(\C[F,Y]u\), direct commutation gives
\begin{equation}\label{eq:two-variable}
\begin{aligned}
E(F^jY^\ell u)
&=j(\mu-j+1-2\ell)F^{j-1}Y^\ell u\\
&\quad+\ell\eta F^jY^{\ell-1}u.
\end{aligned}
\end{equation}
Indeed, \([E,Y]=h_{d_*}\), the \(H\)-weight of \(Y^\ell u\) is
\(\mu-2\ell\), and the further commutators introduce
\(f_{b+d_*}\), which annihilates \(u\).

Define
\[
v_k=\sum_{j=0}^k a_jF^jY^{k-j}u,\qquad a_0=1,\qquad
a_j=-\frac{(k-j+1)\eta}{j(\mu-2k+j+1)}a_{j-1}.
\]
The denominators are nonzero because
\(\mu-2k+j+1\le\mu-k+1<0\).
In \eqref{eq:two-variable}, the coefficient of
\(F^{j-1}Y^{k-j}u\) in \(Ev_k\) is
\[
j(\mu-2k+j+1)a_j+(k-j+1)\eta a_{j-1}=0.
\]
Thus \(Ev_k=0\) and \(Hv_k=\nu v_k\), where \(\nu=\mu-2k\le-2\).
For every \(m\ge1\), in the full localization,
\[
E(F^{-m}v_k)=-m(\nu+m+1)F^{-m-1}v_k.
\]
At \(m=-\nu-1\) the coefficient is zero, and hence
\[
E^{-\nu-1}(F^{-1}v_k)=0.
\]
Passing to the quotient gives
\[
F^{-1}v_k+M=Y^kw_1\ne0.
\]
All terms with \(j\ge1\) lie in \(M\) after multiplication by \(F^{-1}\);
the remaining term is a nonzero localized PBW basis vector.
This proves the claimed nonzero element of \(G_c\).
If \(d_*>L\), then \(\eta=0\) and the construction simplifies to
\(v_k=Y^ku\).
\end{proof}

\begin{theorem}[Deep-region simplicity and the simple head]\label{thm:deep}
Assume \eqref{eq:affine-assumptions}. Then
\[
\boxed{\quad
\T_{f_c}M_{a,b}(\chi)\text{ is simple}
\quad\Longleftrightarrow\quad
\lambda_0-c\kappa\notin\Z.
\quad}
\]
For every value of \(\mu=\lambda_0-c\kappa\), \(G_c\) is the unique
maximal submodule of \(Q_c\), \(Q_c/G_c\) is simple, and \(Q_c\)
is indecomposable. When \(\mu\in\Z\), one has \(0\ne G_c\ne Q_c\),
every proper submodule is contained in \(G_c\), and
\begin{equation}\label{eq:nonsplit}
0\longrightarrow G_c\longrightarrow Q_c\longrightarrow Q_c/G_c
\longrightarrow0
\end{equation}
does not split.
\end{theorem}
\begin{proof}
The FGXZ criterion supplies simplicity of \(M\); PBW supplies
injectivity of \(F\) and semisimplicity of \(H\); the beginning of
this section supplies local nilpotence of \(E\).
Apply Theorem~\ref{thm:containment}, and then
Proposition~\ref{prop:G}.
If \(\mu\notin\Z\), every proper submodule is contained in
\(G_c=0\), so \(Q_c\) is simple.
If \(\mu\in\Z\), the nonzero proper submodule \(G_c\) rules out
simplicity, and the remaining assertions follow from the same theorem.
\end{proof}

\begin{remark}[Where simplicity of the original module enters]
The only use of simplicity of \(M\) in the new submodule argument is
the implication
\[
W(P)\cap M\ne0\quad\Longrightarrow\quad M\subseteq W(P).
\]
The pure-fraction recurrence and the construction of \(G_c\) are
valid without \(\lambda_L\ne0\) or \(\kappa\ne-2\).
Theorem~\ref{thm:deep}, however, is stated for the simple original
modules and must not be asserted for the raw critical or degenerate
induced modules on the basis of this proof.
\end{remark}

\section{The basic case and the extent of the integral classification}

\begin{corollary}\label{cor:basic}
For \(M=M_{-1,1}(\chi)\), assume
\(\lambda_1\ne0\) and \(\kappa\ne-2\). Then
\[
\boxed{\quad
\T_{f_0}M_{-1,1}(\chi)\text{ is simple}
\quad\Longleftrightarrow\quad \lambda_0\notin\Z.
\quad}
\]
For \(\lambda_0\in\Z\), the submodule
\[
G_0=\{v:e_0^nv=0\text{ for some }n\ge1\}
\]
is the unique maximal submodule, \(Q_0/G_0\) is simple, \(Q_0\) is
indecomposable, and its extension by \(G_0\) is nonsplit.
\end{corollary}
\begin{proof}
Here \(L=d_*=1\), \(E=e_0\), \(F=f_0\), \(H=h_0\), and
\(\mu=\lambda_0\), so Theorem~\ref{thm:deep} applies.
\end{proof}

The cyclicity statements in the unified note remain unchanged:
\[
Q_0=Uw_1\quad\Longleftrightarrow\quad
\lambda_0\notin\{-2,-3,\ldots\};
\]
if \(\lambda_0=-s\) with \(s\ge2\), then \(Q_0=Uw_s\).
Indeed, the recurrence
\[
e_0w_m=-m(\lambda_0+m+1)w_{m+1}
\]
produces the pure fractions in the indicated ranges, and
\(f_0w_m=w_{m-1}\) supplies the lower ones.
To pass from pure fractions to arbitrary localized PBW vectors,
give numerator words degree \(2\#e+\#h\).
Commuting a word \(X\) past \(F^{-m}\) gives
\[
Xw_m=F^{-m}Xu+M+
\text{terms of strictly lower numerator degree}.
\]
Induction on this degree, simultaneously for all \(m\), shows that
all pure fractions generate the whole quotient.
When \(\lambda_0=-s\), \(w_1\in G_0\) but \(w_s\notin G_0\), so
\(Uw_1\) is proper.

Cyclicity and simplicity therefore have different exceptional sets.
For example, at \(\lambda_0=0\) the vector \(w_1\) generates \(Q_0\),
but \(Q_0\) is reducible, with
\[
0\ne f_1^2w_1\in G_0,\qquad e_0^3f_1^2w_1=0.
\]

For integral \(\mu\), the additional progress determines the
\emph{head} \(Q_c/G_c\), meaning its maximal semisimple quotient,
which here is simple. It does not determine the internal structure
of \(G_c\). In particular, the following remain open in this note:
\begin{enumerate}[label=\textup{(\roman*)}]
\item whether \(G_c\) is simple;
\item whether \(G_c\) is the socle of \(Q_c\), the sum of its simple
submodules;
\item whether \(Q_c\) has finite composition length, and specifically
whether that length is two.
\end{enumerate}
The existence of a unique maximal submodule must not be confused
with any of these stronger assertions.

\section{Central reduction commutes with quotient localization}

The results of this section do not require the simplicity or deep-region
hypotheses \eqref{eq:affine-assumptions}.

\begin{proposition}[Algebraic reduction by commuting endomorphisms]
\label{prop:reduction}
Let \(V\) be a \(U\)-module, let \(U_F\) be an Ore localization flat
as a right \(U\)-module, and let a commutative algebra \(R\) act on
\(V\) by \(U\)-module endomorphisms. This action extends naturally
to \(D_FV\) and descends to \(\T_FV\).
For any algebraic ideal \(I\subseteq R\), there are natural
isomorphisms
\[
D_F(V/IV)\cong D_FV/I(D_FV),
\qquad
\boxed{\quad
\T_F(V/IV)\cong(\T_FV)/I(\T_FV).
\quad}
\]
In particular, injectivity of \(F\) on \(V/IV\) is not required.
\end{proposition}
\begin{proof}
For \(z\in R\), the extension is
\[
z(s\otimes v)=s\otimes zv\qquad(s\in U_F,\ v\in V).
\]
It is well defined because \(z\) commutes with the \(U\)-action.
The localization map \(\ell:V\to D_FV\) is \(R\)-linear.

Localizing \(0\to IV\to V\to V/IV\to0\) and using flatness yields
\[
D_F(V/IV)\cong D_FV/\operatorname{im}D_F(IV).
\]
The image in the denominator equals \(I(D_FV)\): for instance,
\(F^{-m}(zv)=z(F^{-m}v)\), and every vector in an algebraic ideal
multiple is a finite sum of such terms.
This proves the first isomorphism.

The image of \(V/IV\) in \(D_FV/I(D_FV)\) is
\[
\bigl(\ell(V)+I(D_FV)\bigr)/I(D_FV).
\]
Its cokernel is therefore
\[
\frac{D_FV}{\ell(V)+I(D_FV)}.
\]
On the other hand, \(\T_FV=D_FV/\ell(V)\), and quotienting this
module by \(I(\T_FV)\) gives precisely the same expression.
No embedding of \(V/IV\) into its localization was used.
\end{proof}

\subsection{The actual Sugawara action on the localization}

A \(\g\)-module \(V\) is \emph{smooth} if, for each \(v\in V\),
there is \(N\) with
\[
(\slc\otimes t^N\C[t])v=0.
\]
For \(F=f_c\), localization preserves smoothness.
Indeed, for \(m\ge1\),
\begin{align}
h_rF^{-m}
&=F^{-m}h_r+2mF^{-m-1}f_{r+c},\label{eq:smooth-h}\\
e_rF^{-m}
&=F^{-m}e_r
-mF^{-m-1}(h_{r+c}+r\delta_{r+c,0}K)\notag\\
&\quad-m(m+1)F^{-m-2}f_{r+2c}.\label{eq:smooth-e}
\end{align}
The \(f_r\)'s commute with \(F\). For a fixed numerator \(v\), choose
\(r,r+c,r+2c\) sufficiently large that every resulting current
annihilates \(v\) and the central delta vanishes.
The formulas apply to the image of \(V\) in \(D_FV\) even when the
localization map is not injective.

At level \(K=-2\), the Segal--Sugawara modes are the normally ordered
operators
\begin{equation}\label{eq:sugawara}
T(n)=\frac14\sum_{j\in\Z}
\bigl(2:e_{-j}f_{n+j}:
     +2:f_{-j}e_{n+j}:
     +:h_{-j}h_{n+j}:\bigr).
\end{equation}
For current modes, the convention is
\[
:x_p y_q:=
\begin{cases}
x_p y_q,&p<0,\\
y_q x_p,&p\ge0.
\end{cases}
\]
The sums in \eqref{eq:sugawara} are finite on each vector of a smooth
module. At critical level they commute with the \(\g\)-action
\cite[Section 2.5]{FGXZ}.
They are operators from a completion, not ordinary finite sums in \(U\).

\begin{lemma}[Agreement of the extended and completed actions]
\label{lem:completed-action}
Let \(V\) be a smooth module of level \(-2\) and \(F=f_c\).
The natural extension of the \(U\)-endomorphism \(T(n)\) from \(V\)
to \(D_FV\) agrees with its actual Sugawara action there. In particular,
\[
T(n)(F^{-m}v)=F^{-m}T(n)v.
\]
\end{lemma}
\begin{proof}
The module \(D_FV\) is smooth by
\eqref{eq:smooth-h}--\eqref{eq:smooth-e}, so its Sugawara action is
defined. The map \(V\to D_FV\) intertwines \(T(n)\): it intertwines
each finite current product, and normal ordering makes the relevant
sums finite on each vector.
On \(D_FV\), \(T(n)\) commutes with \(F\).
Since \(F\) is invertible there, it also commutes with \(F^{-1}\).
This gives the displayed formula and determines its value on every
localized vector, exactly as in Proposition~\ref{prop:reduction}.
\end{proof}

\begin{corollary}[Critical central reduction, at any root-mode index]
\label{cor:critical-reduction}
Let \(V\) be a smooth module of level \(-2\). Fix integers \(n_i\)
and scalars \(\theta_i\), indexed by any set \(J\), and let
\[
R=\C[z_i:i\in J],\qquad z_i\text{ act as }T(n_i),\qquad
I=(z_i-\theta_i:i\in J).
\]
For every \(c\in\Z\),
\[
\boxed{\quad
\T_{f_c}(V/IV)\cong(\T_{f_c}V)/I(\T_{f_c}V).
\quad}
\]
The same assertion holds for any algebraic ideal in \(R\).
\end{corollary}
\begin{proof}
Apply Proposition~\ref{prop:reduction} and
Lemma~\ref{lem:completed-action}.
\end{proof}

For the critical modules in the unified note, take
\[
V=M_{a,b}(\chi),\qquad \kappa=-2,\qquad
n_i=L-i,\qquad i\in\N.
\]
Then the corollary supplies the missing compatibility justification
for the displayed central-reduction formulas, including nonboundary
indices \(c<b\).
Together with the previously proved boundary isomorphism and critical
spectral-flow invariance of \(T(n)\), it recovers the unchanged
\(\boldsymbol\theta\)-parameters in the boundary flow.
Those boundary results are recalled from \cite{Unified}, not reproved
here.

\begin{remark}[What is not asserted]
The corollary is an isomorphism of cokernels; it does not by itself
prove that either side is nonzero, that \(f_c\) is injective on the
centrally reduced module, or that the resulting quotient is simple.
Also, ideal multiplication here uses finite algebraic sums.
No assertion is made about arbitrary topologically closed ideals,
completed tensor products, or additional completion procedures.
The ordinary-enveloping-algebra annihilator equality from the earlier
note is not used in this argument.
\end{remark}

\section{Scope, remaining questions, and verification status}

Relative to the open problems in \cite{Unified}, the progress is:
\begin{enumerate}[leftmargin=*,label=\arabic*.]
\item \textbf{Basic and deep-region nonintegral simplicity.}
The sufficient directions are now supplied by
Theorem~\ref{thm:containment}. Combined with the explicit integral
obstruction, this gives the exact criteria in
Theorem~\ref{thm:deep} and Corollary~\ref{cor:basic}.
\item \textbf{Integral-parameter structure.}
The unique maximal submodule, the simple head, indecomposability,
and nonsplitting are determined. Simplicity of \(G_c\), the socle,
and composition length remain unresolved.
\item \textbf{Critical central reduction.}
Compatibility with quotient localization is proved for the algebraic
ideals generated by Sugawara operators, at any root-mode index.
This does not give a classification of the resulting simple modules.
\end{enumerate}

In the \emph{shallow region}
\[
-a\le c<b,\qquad\text{equivalently }1\le b-c<L,
\]
the opposite root \(E=e_{-c}\) is a free PBW mode of \(M\);
it is not locally nilpotent on \(M\).
Thus the hypotheses behind Lemma~\ref{lem:casimir} are unavailable.
Injectivity of \(E\) on the quotient, proved in the unified note,
does not replace this missing hypothesis.
Shallow-region simplicity is not settled by the present argument.
Full parameter isomorphism and multiple-mode simplicity classifications
also remain open.

The new simplicity statements concern the algebra without the degree
derivation. If
\(\widetilde{\g}=\g\rtimes\C d\), with \([d,x_r]=rx_r\), then induction
\[
\I(V)=U(\widetilde{\g})\otimes_U V
\]
is exact and commutes with the localization quotients discussed earlier.
It does not generally preserve simplicity. In particular, the raw
induced modules at noncritical level admit further Casimir reductions.
Neither the simple-head conclusion nor indecomposability is being
transferred to \(\I(Q_c)\) without a separate argument.

The definitions and integer-parameter construction recalled above come
from the unified note; the FGXZ simplicity criterion is an external
input. The new proofs are the rank-one containment argument and the
central-reduction compatibility argument.
In the preceding review, the existing Takiff script passed 5,866 exact
finite checks, and the rank-one factorization was also checked by
exact polynomial expansion for the first twelve positive exponents.
These checks are corroboration only; the proofs above apply to all
parameters and all exponents under their stated hypotheses.
No independent peer review, Danus certification of these new results,
or Lean/Coq formal verification is claimed.
This addendum is supplied as LaTeX source without compilation.

\begingroup
\raggedright
\begin{thebibliography}{9}
\bibitem{FGXZ}
V.~Futorny, X.~Guo, Y.~Xue, and K.~Zhao,
\emph{Smooth representations of affine Kac--Moody algebras},
arXiv:2404.03855v2 (2024).
\href{https://arxiv.org/html/2404.03855v2}{Full text}.

\bibitem{Unified}
Project working note,
\emph{Boundary and Nonboundary Localization Quotients of Smooth
Affine Modules}, September 2026.
Source: \texttt{localization\_quotients\_unified\_en.tex}.
\end{thebibliography}
\endgroup
\end{document}
